\documentclass{article}
\usepackage{tikz}
\usepackage{graphicx}
\usepackage[brazilian]{babel}
\usepackage[utf8]{inputenc}
\usepackage[T1]{fontenc}

\begin{document}
	
	\begin{center}
		\begin{tikzpicture}[
			box/.style={draw, rectangle, minimum width=2cm, minimum height=1cm, align=center},
			title/.style={font=\large\bfseries}
			]
			
			% Título
			\node[title] at (0, 8) {Programa Genes};
			
			% Imagem central (placeholder - substitua pelo seu gráfico/logo)
			\node at (0, 4) {\includegraphics[width=3cm]{example-image}};
			% Para usar sua própria imagem, substitua "example-image" pelo nome do arquivo
			% Exemplo: \includegraphics[width=3cm]{logo.png}
			
			% Caixas do lado ESQUERDO (4 caixas)
			\node[box] at (-5, 6) {Nome 1};
			\node[box] at (-5, 4) {Nome 2};
			\node[box] at (-5, 2) {Nome 3};
			\node[box] at (-5, 0) {Nome 4};
			
			% Caixas do lado DIREITO (4 caixas)
			\node[box] at (5, 6) {Nome 5};
			\node[box] at (5, 4) {Nome 6};
			\node[box] at (5, 2) {Nome 7};
			\node[box] at (5, 0) {Nome 8};
			
			% Caixa ACIMA
			\node[box] at (0, 7) {Nome 9};
			
			% Caixa ABAIXO
			\node[box] at (0, 1) {Nome 10};
			
			% Linhas de conexão (opcional - descomente se quiser)
			% \draw[->] (-3.5,6) -- (-1,5);
			% \draw[->] (-3.5,4) -- (-1,4);
			% \draw[->] (-3.5,2) -- (-1,3);
			% \draw[->] (1,5) -- (3.5,6);
			% \draw[->] (1,4) -- (3.5,4);
			% \draw[->] (1,3) -- (3.5,2);
			% \draw[->] (0,6.5) -- (0,5);
			% \draw[->] (0,2) -- (0,3);
			
		\end{tikzpicture}
	\end{center}
	
\end{document}